


\section{A quantitative boundary argument}
Default, covenant breach and a margin call are threshold events. Let $X$ be
exact cash at a decision date, $Y$ its approximation under a specified joint
coupling, and $c$ the contractual threshold. All three are in the same currency
and date basis. A cash approximation of $-\epsilon$ by $+\epsilon$ has error
$2\epsilon$, yet changes $\ind\{X<0\}$ from one to zero. Small cash error
alone does not imply small default-probability error.

\begin{proposition}
If $\E|X-Y|<\infty$, then for every $\delta>0$,
\begin{equation}\label{sim:lit-threshold}
|\Prob(X<c)-\Prob(Y<c)|
\leq \Prob(|X-c|\leq\delta)+\frac{\E|X-Y|}{\delta}.
\end{equation}
\end{proposition}
\begin{proof}
The absolute difference in event probabilities is at most the probability
that their indicators disagree. If $|X-c|>\delta$ and $|X-Y|<\delta$, $X$
and $Y$ lie on the same side of $c$. Consequently disagreement is contained
in $\{|X-c|\leq\delta\}\cup\{|X-Y|\geq\delta\}$. The union bound
and Markov's inequality give (\ref{sim:lit-threshold}).
\end{proof}
The first term is probability mass close to the covenant; the second is the
probability allowed for a sufficiently large approximation error. This local
derivation is independent of any claim that a particular algorithm supplies
the required coupling error. Write $e=\E|X-Y|$. If $e=0$, then $X=Y$ almost surely and the threshold probabilities agree.
For $e>0$, if $X$ has a density bounded by $M>0$ near $c$,
the first term is at most $2M\delta$. Choosing
$\delta=\sqrt{e/(2M)}$, when that choice lies within the neighborhood where
the density bound holds, gives $2\sqrt{2Me}$. Atoms at the threshold and
uncontrolled conditional distributions prevent this simplification.

For several dates, a union bound can bound the probability that any breach
indicator differs by the sum of the per-date bounds, provided the relevant
coupling and cash errors are supplied. If a different breach decision changes
subsequent actions and cash flows, those cash errors must already account for
that feedback; an open-loop cash comparison does not establish them. When
there are no weights over histories, the appropriate output is instead an
uncertain boundary region or certified feasible-set containment. A probability
bound cannot be manufactured from the proportion of enumerated witnesses.

\section{A common-policy regret argument}
Let $J(\pi)$ and $\widetilde J(\pi)$ be expected costs in two models, with
the same admissible policy class $\Pi$, finite expected costs and finite optimal values. This comparison presupposes that a
policy can be implemented with the information and feasible actions in both
models. Suppose a uniform bound
$\sup_{\pi\in\Pi}|J(\pi)-\widetilde J(\pi)|\leq\varepsilon$ holds.
If $\widehat\pi$ is $\eta$-optimal for the approximate problem, then
\begin{align}
J(\widehat\pi)&\leq\widetilde J(\widehat\pi)+\varepsilon\nonumber\\
&\leq\inf_{\pi\in\Pi}\widetilde J(\pi)+\eta+\varepsilon\nonumber\\
&\leq\inf_{\pi\in\Pi}J(\pi)+\eta+2\varepsilon.
\end{align}
The last step follows by applying the uniform bound to an arbitrarily
near-optimal policy of the exact problem, so attainment is unnecessary.
The result is a regret bound for a common feasible policy. If one only knows
that the two infima are close, the first inequality is missing. If a policy
lift changes the policy class, a bound for that lift and the relevant comparison
class is needed. This is why reporting a small optimized residual or close
portfolio value alone does not establish that a compressed financing policy
will work in the original company.

\section{When sampled paths disclose their own futures}
With continuous first-period observations, independent sampled paths typically
have distinct first values. If each stored first value has only its single
sampled continuation, that value becomes an identifier for the sampled future.
An optimizer can exploit this unintended information. More paths do not by
themselves repair the information structure of the empirical measure.
Backhoff and coauthors give an explicit optimal-stopping counterexample
\cite[Example~7.1]{litcloseestimation}. Their adapted empirical measure quantizes
observations so that paths share histories. Strong consistency applies to
iid complete paths on a fixed finite horizon and bounded state space;
quantitative rates additionally require Lipschitz conditional kernels
\cite[Definitions~1.1--1.2, Theorems~1.3, 1.5, 1.7]{litcloseestimation}.

A Markov modification improves the horizon exponent only when the process
is in fact Markov \cite[Theorem~6.1]{litcloseestimation}; constants still depend on
the horizon. A company's current balance sheet generally is not a sufficient
Markov state when hidden debt terms or asset vintages matter. Nor are reports
from one issuer an iid sample of complete company paths. The method suggests
how a deliberately generated simulation sample can retain information. Its
rates cannot simply be assigned to an observed issuer panel.

\section{A finite coupling method for discrete decisions}
A stronger but directly checkable condition can avoid averaging actions.
Consider a finite horizon of $T$ transitions. Keep a persistent contract type
$\theta$ fixed and let a history record the actor's observations and past
actions. In both the full model and a reduced model, require the same feasible
actions at paired histories. Couple their initial states so that the required
observations agree. At each paired history and common action, suppose a joint
coupling of the next state, observation and reported output leaves the states
paired and the observation/output records equal with conditional probability
at least $1-\delta_t$. This condition includes every history and action reachable
before a mismatch; fitting one historical policy is insufficient.

If $E_t$ is the event that no mismatch has occurred through transition $t$,
then $\Prob(E_t\mid E_{t-1})\ge1-\delta_t$. Repeated conditioning gives
\begin{equation}
\Prob(E_T^c)\le1-\prod_{t=1}^T(1-\delta_t)
\le\sum_{t=1}^T\delta_t.
\label{sim:lit-finite-coupling}
\end{equation}
A common deterministic policy chooses the same action while the observation
histories agree. Randomized policies can use the same independent random draws,
if both models permit that randomization. No convexity or fractional issuance
is required. A bounded path loss $L\in[0,B]$ agrees on $E_T$, and hence
\begin{equation}
|\E L-\widetilde\E L|
\le B\Prob(E_T^c)
\le B\left\{1-\prod_{t=1}^T(1-\delta_t)\right\}.
\label{sim:lit-finite-cost}
\end{equation}
A default indicator is covered because it is a bounded output and the
coupling condition compares its probability explicitly. Small cash distance
alone does not provide that condition. If exact output matching is impossible,
a separate bound on the loss difference on matched paths must be added.

For finite kernels with common outcome labels, the smallest one-step mismatch
probability is their total-variation distance
$\frac12\sum_z|K(z)-\widetilde K(z)|$. Pairing full states with reduced states
requires a coupling whose successful pairs remain in the declared relation;
marginal observation agreement is insufficient. An implementable finite
procedure first partitions states by persistent type, available information,
feasible actions and required outputs. It then compares joint successor
kernels into the current blocks for every action, splitting a block whenever
these differ. Repeat until no block changes. Each split reduces the number of
states per block, so a finite partition stabilizes after finitely many splits.
This yields an exact quotient for the declared finite model, not a claim that
the model contains every economically possible company.

Approximate sharing uses the same comparison with an explicit error allowance
and a checked pairing. Threshold regimes and different feasible actions are
kept separate unless their mismatch has been bounded. An unresolved comparison
is retained as unresolved, not assigned zero error. For a fixed common policy
class, the uniform cost bound in \eqref{sim:lit-finite-cost} also gives the common-policy
$2\varepsilon+\eta$ regret bound derived in this chapter. If the persistent types are unweighted,
apply the bound within each type. A uniform type-wise bound also bounds the
worst-type value, since
$|\sup_\theta J_\theta-\sup_\theta\widetilde J_\theta|
\le\sup_\theta|J_\theta-\widetilde J_\theta|$.
The adversary does not redraw the contract type at each date.

The exact reference uses three periods, cash observed by management, a hidden
fixed type, binary borrowing, and a default penalty. It compares complete
paths for all 64 policies in the explicitly restricted class of date-and-
zero-cash rules. Removing an irrelevant tag preserves every path. Merging
types is rejected. A synthetic change of $1/20$ in each period's bad-expense
probability gives the conservative cost bound
$12\{1-(19/20)^3\}=3423/2000$. The largest observed value discrepancies are
$4433/8000$ and $1/2$ in the two types, within the bound. These finite checks
support the stated method and expose its conservatism. They do not supply a
capacity estimate or a bound for an unenumerated continuous company model.
